\section{Introducción}

El control de la calidad del aire interior en laboratorios constituye un requisito crítico para garantizar la integridad de los procesos, la seguridad del personal y la confiabilidad de los resultados analíticos. En instalaciones sensibles a contaminación por partículas o microorganismos, la ventilación debe cumplir simultáneamente dos funciones: diluir contaminantes generados en el interior y evitar la infiltración de agentes externos. La presurización positiva respecto a zonas adyacentes constituye la estrategia más eficaz para alcanzar el segundo objetivo.

El presente informe aborda el diseño conceptual del sistema de ventilación y presurización de un laboratorio con volumen efectivo de 320 m$^3$, ubicado en las instalaciones de BRINSA. El recinto requiere 12 renovaciones de aire por hora, valor seleccionado con base en los estándares internacionales aplicables a laboratorios de control medio y bioseguridad nivel 2. El aire de impulsión se introduce mediante un ventilador directo, sin ductos de distribución, y sale del recinto a través de tres rejillas de exfiltración controladas. Esta configuración minimiza la complejidad constructiva y facilita la modelación posterior mediante dinámica de fluidos computacional (CFD).

El documento se estructura de la siguiente manera. La Sección 2 presenta los objetivos y el alcance del estudio. La Sección 3 define las bases de diseño, incluyendo condiciones ambientales, criterios de ventilación y normativa aplicable. La Sección 4 describe la metodología de cálculo. La Sección 5 presenta los resultados de caudal, potencia, áreas de impulsión y exfiltración, y los datos para CFD. La Sección 6 analiza la coherencia de los resultados con los requerimientos de presurización. Finalmente, las Secciones 7 y 8 presentan conclusiones y recomendaciones.
